\section{Literature Review}

Deep learning research advances have led to significant interest in applying convolutional neural networks (CNNs) to a wide range of computer vision problems, including breast cancer detection from mammography images (Litjens et al., 2017; Shen et al., 2017). CNN-based models learn hierarchical feature representations directly from pixel-level data, enabling the detection of subtle visual patterns that may be difficult to identify using traditional rule-based or handcrafted feature approaches. In contrast, earlier computer-aided detection (CAD) systems relied heavily on predefined features and heuristic rules, which limited their ability to generalise across diverse patient populations and imaging conditions, particularly in cases involving dense breast tissue (Elmore et al., 2015; Dromain et al., 2013).

As of 2025, CNNs appear to have become the dominant modelling approach in mammography research due to their ability to automatically extract discriminative spatial features from medical images. Numerous studies report strong performance when applying CNNs to classify mammograms as benign or malignant under controlled experimental conditions (Shen et al., 2017; McKinney et al., 2020; Ragab et al., 2022).

In the review published by Wang et al. (2024) they examined over 60 deep learning-based mammography studies and reported area under the curve (AUC) values ranging between 0.77 to 0.98. AUC measures how well a model can distinguish between classes by summarising its performance across all classification thresholds; a higher AUC indicates better overall discriminative ability. While these results demonstrate strong discriminative performance under controlled conditions, the review highlights substantial variation in evaluation methodology, including inconsistent use of metrics and frequent reliance on image level rather than patient level validation. This raises concerns regarding data leakage and inflated performance estimates, particularly in screening contexts where generalisation to unseen patients is critical. This data leakage can occur when information from outside the training data is unintentionally used to build a model, leading to unrealistically good performance and poor generalisation to new data.

The availability of publicly accessible datasets has played a crucial role in aiding deep learning research in mammography. The Digital Database for Screening Mammography (DDSM) has generally been one of the most widely used datasets; however, inconsistencies in annotation and metadata posed challenges for reproducibility (Heath et al., 2001; Lee et al., 2017).

To address these limitations, Lee et al. (2017) introduced the Curated Breast Imaging Subset of DDSM (CBIS-DDSM), which provides standardised labels, preprocessing pipelines, and patient identifiers. This dataset has since become a benchmark for CNN-based mammography studies and supports more rigorous experimental design, including patient level splitting strategies.

One issue highlighted within the literature review is data scarcity; there are too few high quality, labelled examples available to train a model effectively; CBIS-DDSM is still relatively small when compared to datasets typically used to train deep neural networks in other domain areas. This data scarcity presents challenges for training deep architectures from scratch, increasing the risk of overfitting.

To address this, numerous studies have adopted transfer learning approaches, leveraging pre-trained networks such as VGC, ResNet and EfficientNet.(Shen et al., 2017; Ragab et al., 2022; Huynh et al., 2016) Transfer learning models are pre-trained on large scale natural image datasets and are fine-tuned for domain specific tasks.

Ragab et al. (2022) reviewed the application of transfer learning in breast cancer detection and reported that fine-tuning pre-trained CNNs generally improves performance on small mammography datasets compared to training from scratch. This improvement is attributed to the reuse of low-level and mid-level feature representations learned from large datasets, which can accelerate convergence and reduce overfitting.

However, Ragab et al. (2022) also caution that transfer learning introduces potential risks related to domain mismatch, as natural images differ substantially from medical images in texture, scale, and semantic structure. Without careful fine-tuning and regularisation, transferred features may encode biases that reduce generalisability. These concerns highlight the need for careful architectural and training design when applying transfer learning in medical contexts.

In 2025, Kaddes et al. (2025), hybrid architectures have been explored that combine CNN features extraction with sequence-based models. Kaddes et al. (2025) proposed a CNN-LSTM hybrid model, where the CNN extracts spatial features from mammography images and the long short-term memory (LSTM) network models sequential dependencies within the extracted features; it reported an accuracy of 98.4\% on the CBIS-DDM dataset, demonstrating the potential of hybrid architectures to improve predictive performance. While these results illustrate the potential of advanced architectures, accuracy is seen as the primary metric and they do not evaluate performance at clinically relevant operating points, which would be something clinicians would need data on to ascertain if these models would be fit for purpose in real-world medical scenarios and be trusted by medical professionals that would seek to use them to aid in their work.

Across the reviewed literature, three modelling strategies appear most prominent: CNNs trained from scratch, transfer learning based CNNs and hybrid architectures (Wang et al., 2024; Ragab et al., 2022). Each approach prioritises different objectives and presents different trade-offs as a result.

CNNs trained from scratch offer greater domain specific feature learning and architectural transparency but are heavily constrained by data scarcity. Transfer learning improves data efficiency and performance on small datasets but introduces reliance on features learned from non-medical domains. Hybrid architectures prioritise raw performance by modelling higher order relationships but increase model complexity and reduce interpretability (Ragab et al., 2022; Wang et al., 2024).

No single approach consistently addresses all challenges associated with mammography screening. Instead, the various papers suggest that methodological rigor in validation and evaluation may be as important as architectural sophistication in achieving clinically meaningful results (Wang et al., 2024; McKinney et al., 2020).

Evaluation strategy is a critical factor in determining the real-world usefulness of deep learning models. Many mammography studies report accuracy or ROC-AUC as primary metrics; however, these measures may be misleading in highly imbalanced datasets, where cancer-positive cases are rare. When considering the number of patients that are screened for breast cancer in the UK alone per year, compared to the number of mammograms that test positive for cancer, the balance skews massively towards negative results. According to NHS Digital, around 2.2 million breast cancer screenings took place in 2022, of those, around 19,000 detected cancer in the patient; around 0.9\%, meaning around 99\% of all screenings return a negative screening result. From a human perspective, this is a good thing as nobody wants to be diagnosed with cancer, however, from a data science perspective, this creates a high data imbalance. 

In such scenarios, high accuracy can be achieved even when a model performs poorly on the minority class. Saito and Rehmsmeier (2015) demonstrated that precision–recall metrics and F1-score provide more informative evaluation in imbalanced classification tasks, as they explicitly capture the balance between false positives and false negatives.

From a clinical perspective, false negatives are particularly critical, as missed diagnoses may delay treatment and worsen patient outcomes. Screening programmes often operate at fixed specificity thresholds, making sensitivity at fixed specificity a more meaningful metric than aggregate AUC. Despite this, Wang et al. (2024) note that few studies report such clinically relevant measures, limiting the interpretability of reported performance.

Of the more recent papers (McKinney et al., 2020; Topol, 2019) emphasis has been increasing in deep learning systems to be designed as decision support tools rather than autonomous diagnostic ages, creating more of a human in the loop AI tool. Generally, clinicians require the highest level of reliability, so the ability for an expert to override autonomous decisions when there is a level of uncertainty is something that should be considered.

Human in the loop approaches, where models flag high risk or uncertain cases for further review would reduce the risk of becoming over reliant on algorithmic output and align with ethical and regulatory requirements (Topol, 2019; EU AI Act, 2023). However, it seems that a principle following a type of uncertainty estimation or a confidence calibration to involve a human seems relatively underexplored from the literature reviewed in this case study.

One of the main challenges highlighted in the literature, and one that would benefit from such a flagging system, is the variation in breast density, which significantly affects both human and algorithmic interpretation of images. Dense breast tissue appears radiopaque on mammograms, often obscuring lesions and reducing contrast between malignant and normal regions. Studies have shown (Boyd et al., 2007; Kerlikowske et al., 2010) that breast density is associated with increased rates of false negatives in traditional screening, highlighting an important limitation of both radiologist led and automated detection approaches.

Deep learning models can partially mitigate this challenge by learning more complex texture patterns beyond simple density intensity differences. However, Wang et al. (2024) note that many studies fail to stratify performance by breast density, making it difficult to assess whether reported improvements generalise across different patient subgroups. Without density aware evaluation, models may demonstrate strong average performance while still underperforming in clinically challenging cases. This omission represents a further gap in the literature, particularly given the clinical importance of dense breast tissue in screening populations, though gathering the data required for such a task presents an additional issue in itself.

Another significant challenge is generlisation beyond the training dataset, which is a challenge for machine learning in all medical settings, not just breast cancer screening. Many mammography studies evaluate models exclusively on internal test splits derived from the same dataset used for training. While such evaluations may indicate strong performance, they provide limited insight into how models will behave when exposed to data from different institutions, imaging devices, or patient demographics. Something that can be an issue as mentioned earlier with the imbalanced datasets in real world scenarios where the proportion of positive tests is under 1\% of the dataset.

Wang et al. (2024) highlight that only a small subset of reviewed studies perform external validation on independent datasets. This lack of external testing raises concerns regarding reliability and real-world applicability. Differences in image acquisition protocols, scanner manufacturers, and population characteristics can substantially affect model performance. For example, from what I was able to find, there were over 25 different manufacturers for mammography scanners, some of which use slightly different techniques. The Siemens Healthineers system generally uses direct conversion detectors, which give a sharp, high frequency noise profiles, whereas the GE HelathCare systems use indirect conversion detectors, which result in images with a smoother texture (Pisano et al., 2008; Vedantham et al., 2015). Consequently, models that perform well on curated benchmark datasets may experience performance degradation when deployed in clinical environments due to some of these hard-to-calculate variations.

Although full external validation may be beyond the scope of my project, being aware of these limitations will help in my evaluation of the success of the final output. 

In a healthcare setting, interpretability is generally recognised as a key requirement when deciding whether to adopt deep learning systems in healthcare, understanding how the decisions have been made is critical. A criticism sometimes directed towards CNNs is that they are ‘block box’ models (Rudin, 2019), producing predictions without explicit explanations of the underlying reasoning. From the point of view of a medical care provider, the lack of transparency can hinder trust and raise concerns.

Other studies (Selvaraju et al., 2017) attempt to address interpretability through techniques such as saliency maps or gradient based visualisations, which highlight image regions contributing most strongly to model predictions. However, Wang et al. (2024) caution that such visualisations are not always reliable indicators of true model reasoning and may provide misleading reassurance if interpreted uncritically.

The literature therefore suggests a trade-off between predictive performance and interpretability. More complex architectures, including hybrid models, may achieve higher accuracy but further obscure the decision-making process, a black box-like scenario. This trade-off reinforces the importance of positioning deep learning systems as decision support tools rather than autonomous diagnostic agents, particularly in high risk screening scenarios.

Beyond technical performance, ethical considerations play a crucial role in evaluating the suitability of AI systems for medical use. Breast cancer screening affects large populations, and errors in automated systems may disproportionately impact vulnerable groups if biases are present in training data. Dataset composition, including demographic representation and disease prevalence, therefore has important ethical implications.

The literature increasingly emphasises the need for transparent reporting of dataset characteristics. However, many mammography studies provide limited demographic detail, making it difficult to assess potential biases.

Ethical deployment also requires consideration of accountability. In human-in-the-loop systems, responsibility for final decisions remains with clinicians, whereas fully automated systems raise complex legal and ethical questions. The prevailing consensus within the literature supports AI-assisted screening frameworks that enhance, rather than replace, clinical expertise.

The practical value of deep learning models depends on accuracy and on how they integrate into existing workflows. Within the literature, there are suggestions of using AI systems as pre-screening tools to prioritise high-risk cases, while others suggest post-screening review to reduce false negatives. Designing models that output confidence scores or uncertainty estimates may support more flexible workflow integration, allowing clinicians to focus attention where algorithmic certainty is lowest.

From the literature that I reviewed, there is a clear demonstration that there has been substantial progress in applying deep learning to breast cancer detection from mammography images. 

CNN-based approaches, transfer learning, and hybrid architectures have all achieved high reported performance on curated datasets. However, closer examination reveals persistent limitations related to validation methodology, evaluation metrics, interpretability, and clinical integration.

Performance metrics frequently prioritise aggregate measures such as accuracy and AUC, despite their limited relevance in imbalanced screening contexts. Image level validation and limited external testing raise concerns regarding data leakage and generalisability. Furthermore, increased model complexity often comes at the expense of interpretability and transparency, which are critical for clinical trust.

Collectively, these findings highlight a gap between research performance and reliability if deployed in a real medical setting. Addressing this gap requires careful consideration of evaluation strategy, validation design, and clinical context. My project will be motivated by these limitations, and I will look to adopt patient level validation, clinically meaningful evaluation metrics, and decision support-oriented modelling to provide a more realistic assessment of deep learning-based mammography detection systems.
